\newcommand{\Date}{March 25, 2019}
\newcommand{\Title}{Bible Notes: Amos 2}

\documentclass[a4paper,10pt]{article}

% Lwarp
%\usepackage[%
    %latexmk,
    %mathjax,
%]{lwarp}

% Bibliography
\usepackage[
    citepages=omit,
    fullbibrefs,
    ibidtracker=false,
    idemtracker=false,
    style=sbl,
]{biblatex}
\addbibresource{../../bib/books.bib}
\addbibresource{../../bib/commentaries.bib}
\addbibresource{../../bib/commentary-series.bib}
\addbibresource{../../bib/dictionaries-encyclopedias.bib}
\addbibresource{../../bib/general-series.bib}
\addbibresource{../../bib/journal-articles.bib}
\addbibresource{../../bib/websites.bib}
\setlength{\bibhang}{0.5in}

% Packages with arguments
\usepackage[bookmarks,hidelinks]{hyperref}
\usepackage[rm,tiny]{titlesec}
\usepackage[english]{babel}
\usepackage[margin=1in]{geometry}
\usepackage[marginal]{footmisc}

% Packages without arguments
\usepackage{afterpage}
\usepackage{booktabs}
\usepackage{fancyhdr}
\usepackage{graphicx}
\usepackage{longtable}
\usepackage{paralist}
\usepackage{ragged2e}
\usepackage{scrextend}
\usepackage{setspace}
\usepackage{times}
\usepackage{tocloft}
\usepackage{xcolor}

% Quotation marks -- British style
\usepackage[autostyle,english=british]{csquotes}
%\DeclareQuoteStyle[american]{english}
    %{\textquotedblleft}
    %[\textquotedblleft]
    %{\textquotedblright}
    %[0.05em]
    %{\textquoteleft}
    %{\textquoteright}
%\DeclareQuoteStyle[american-verse]{english}
    %{\textquotedblleft}
    %{\textquotedblright}
    %[0.05em]
    %{\textquoteleft}
    %{\textquoteright}
\DeclareQuoteStyle[british]{english}
    {\textquoteleft}
    [\textquoteleft]
    {\textquoteright}
    [0.05em]
    {\textquotedblleft}
    {\textquotedblright}
\DeclareQuoteStyle[british-verse]{english}
    {\textquoteleft}
    {\textquoteright}
    [0.05em]
    {\textquotedblleft}
    {\textquotedblright}
% British
\newcommand{\OpenQuote}{\textquoteleft}
% American
%\newcommand{\OpenQuote}{\textquotedblleft}

% Poetry
\newcommand{\VerseQuoteStyle}{\setquotestyle[british-verse]{english}}
\newcommand{\NormalQuoteStyle}{\setquotestyle[british]{english}}
\newcommand{\VerseIndent}{\hspace*{2em}}
\newcommand{\VerseIndentTwo}{\hspace*{4em}}
\newcommand{\VerseIndentFour}{\hspace*{8em}}

% Lengths
\setlength{\parindent}{0.5in}
\setlength{\RaggedRightParindent}{\parindent}

% Line spacing
\doublespacing{}

% Justification
\RaggedRight

% Footnotes
\setlength{\footnotesep}{19.86pt} % Determined by running \footnotesize \the\baselineskip
\renewcommand{\footnoterule}{\noindent\smash{\rule[3pt]{2in}{0.4pt}}\vspace{-0.5\footnotesep}}
\setlength{\footnotemargin}{0.5in}
\deffootnote[0.5in]{0pt}{0.5in}{}
\let\TempFootnote\footnote
\renewcommand{\footnote}[1]{\TempFootnote{\thefootnotemark.\enskip#1}}
\setlength{\skip\footins}{\baselineskip}

% Table of contents
\renewcommand\cftsecfont{\rm}
\renewcommand\cftsecleader{\cftdotfill{.}}
\renewcommand\cftsubsecdotsep{\cftnodots}

% Word count
%TC:newcounter fwords Words in footnotes
%TC:newcounter footnote Number of footnotes
%TC:macro \footnote [fwords]
%TC:macroword \footnote [footnote]
%TC:macro \autocite [fwords]
%TC:macroword \autocite [footnote]

\begin{document}

% Sections
\setcounter{secnumdepth}{0}
% - Section
\titlespacing*{\section}{0pt}{\baselineskip}{\baselineskip}
\titleformat*{\section}{\center\uppercase}
% - Subsection
\titlespacing*{\subsection}{0pt}{\baselineskip}{0pt}
\titleformat*{\subsection}{\center\bfseries}

% Headers & Footers
\renewcommand{\headrulewidth}{0pt}
\fancypagestyle{first}{%
    \fancyhf{}
    \fancyfoot[C]{\thepage}
}
\fancypagestyle{subsequent}{%
    \fancyhf{}
    \fancyhead[R]{\thepage}
}
\pagestyle{empty}
\pagenumbering{roman}

\mbox{}

\vfill

\begin{center}
    \Huge{\textbf{\Title}}
\end{center}


\vfill

\begin{center}
    \Date

    Jason O'Conal
\end{center}

\vfill
\newpage

\pagestyle{plain}

\mbox{}\vspace{1in} % Extra margin

\renewcommand{\contentsname}{}
\section*{CONTENTS}
\tableofcontents

\vfill
\newpage

\mbox{}\vspace{1in} % Extra margin

\printbiblist[title=ABBREVIATIONS]{abbreviations}

\vfill
\newpage

\pagestyle{subsequent}

\thispagestyle{first}

\pagenumbering{arabic}

\mbox{}\vspace{1in} % Extra margin


\section{Introduction}

From 1:3, Amos has been closing in on Israel.

\begin{enumerate}
    \item First, nations from Israel's political environment:
        \begin{itemize}
            \item Syria (1:3--5)
            \item Philistia (1:6--8)
            \item Tyre (1:9--10)
        \end{itemize}
    \item Then, Israel's historical cousins:
        \begin{itemize}
            \item Edom (1:11--12)\footnote{%
                    \enquote{Once when Jacob was cooking some stew, Esau came in
                    from the open country, famished. He said to Jacob,
                    \enquote{Quick, let me have some of that stew! I'm
                    famished!} (That is why he was also called Edom.} (Genesis
                    25:29f)
                }
            \item Ammon (1:13--15))\footnote{%
                    \enquote{So both of Lot's daughters became pregnant by their
                    father. The older daughter had a son, and she named him
                    Moab; he is the father of the Moabites today. The younger
                    daughter also had a son, and she named him Ben-Ammi; he is
                    the father of the Ammonites of today.} (Genesis 19:36--38)
                }
            \item Moab (2:1--3)
        \end{itemize}
    \item Then, Israel's brother, Judah (2:4--5)
    \item Finally, Israel herself (2:6--16)\footnote{%
            \enquote{The first three nations reviewed, Syria, Philistia and
            Tyre, belonged to Israel's political environment and nothing more.
            But in the biblical view of things the next three, Edom
            (\textit{cf.} Gn.  25:29f), Ammon and Moab (\textit{cf.} Gn.
            19:36--38), were \enquote{cousins} to Israel: divine judgment was
            now falling, so to say, within the family.}
            \autocite[50]{motyer:2011}}
\end{enumerate}

\begin{center}
    \includegraphics[width=0.5\textwidth]{amos-2-target.eps}
\end{center}

I looked at this briefly in my post on Amos 1, but now that we turn our focus to
Judah ans Israel, it's worth repeating: the \enquote{ground of accusation is
never simply that they have acted to the detriment of Israel}\footnote{%
    \enquote{It is appropriate to notice here that in condemning these
    nations---both the outright heathen and the related heathen---the ground of
    accusation is never simply that they have acted to the detriment of Israel.}
    \autocite[50]{motyer:2011}
}. This is important because Judah and Israel seem to
have taken God's choice of them as his people, God's choice to make his dwelling
in Jerusalem, as a licence to do whatever they felt like.

In this post, I'll be looking at the oracles against Judah and Israel in Amos
2:4--16. I'll focus on what these oracles have to say to Christians
today.\footnote{%
    I will not engage with questions of the source and unity of Amos in this
    post. I may look at it in future posts, though I make no promises. For those
    who would like to read about this, Smith's \enquote{Unity} section is well
    written and provides a good starting point.
    \autocite[See][110--114]{smith:2017}
}

\section{Amos 2:4--16}

\VerseQuoteStyle
\begin{quote}
    This is what the \textsc{Lord} says:

    \VerseIndent\enquote{For three sins of Judah,\\
    \VerseIndentTwo even for four, I will not relent.\\
    \VerseIndent Because they have rejected the law of the \textsc{Lord}\\
    \VerseIndentTwo and have not kept his decrees,\\
    \VerseIndent because they have been led astray by false gods,\\
    \VerseIndentTwo the gods their ancestors followed,\\
    \VerseIndent I will send fire on Judah\\
    \VerseIndentTwo that will consume the fortresses of Jerusalem.}

    This is what the \textsc{Lord} says:

    \VerseIndent\enquote{For three sins of Israel,\\
    \VerseIndentTwo even for four, I will not relent.\\
    \VerseIndent They sell the innocent for silver,\\
    \VerseIndentTwo and the needy for a pair of sandals.\\
    \VerseIndent They trample on the heads of the poor\\
    \VerseIndentTwo as on the dust of the ground\\
    \VerseIndentTwo and deny justice to the oppressed.\\
    \VerseIndent Father and son use the same girl\\
    \VerseIndentTwo and so profane my holy name.\\
    \VerseIndent They lie down beside every altar\\
    \VerseIndentTwo on garments taken in pledge.\\
    \VerseIndent In the house of their god\\
    \VerseIndentTwo they drink wine taken as fines.

    \VerseIndent\OpenQuote Yet I destroyed the Amorites before them,\\
    \VerseIndentTwo though they were tall as the cedars\\
    \VerseIndentTwo and strong as the oaks.\\
    \VerseIndent I destroyed their fruits above\\
    \VerseIndentTwo and their roots below.\\
    \VerseIndent I brought you out of Egypt\\
    \VerseIndentTwo and led you forty years in the wilderness\\
    \VerseIndentTwo to give you the land of the Amorites.

    \VerseIndent\OpenQuote I also raised up prophets from among your children\\
    \VerseIndentTwo and Nazirites from among your youths.\\
    \VerseIndent Is this not true, people of Israel?}\\
    \VerseIndentFour declares the \textsc{Lord}.\\
    \VerseIndent\enquote{But you made the Nazirites drink wine\\
    \VerseIndentTwo and commanded the prophets not to prophesy.

    \VerseIndent\OpenQuote Now then, I will crush you\\
    \VerseIndentTwo as a cart crushes when loaded with grain.\\
    \VerseIndent The swift will not escape,\\
    \VerseIndentTwo the strong will not muster their strength,\\
    \VerseIndentTwo and the warrior will not save his life.\\
    \VerseIndent The archer will not stand his ground,\\
    \VerseIndentTwo the fleet-footed soldier will not get away,\\
    \VerseIndentTwo and the horseman will not save his life.\\
    \VerseIndent Even the bravest warriors\\
    \VerseIndentTwo will flee naked on that day,}\\
    \VerseIndentFour declares the \textsc{Lord}.
\end{quote}
\NormalQuoteStyle

\section{Judah and Israel}

\subsection{The Oracle Against Judah}

The preceding oracles (1:3--2:3) reference crimes of conscience. There were no
written rules that the nations were breaking, the focus was instead on the
unwritten rules written on people's consciences. The Judahites had the same
requirements and may have been guilty of similar crimes but they also had God's
law and so they would be judged by that as well. My focus in this post is
mainly the oracle against Judah, but I think the other oracles set the scene for
the oracle against Israel: they gradually gain more and more acceptance of the
standards that God uses to judge: the Israelites and the Judahites had the same
God-given consciences as the other nations; the Israelites had the same law as
Judah.

\begin{quote}
    On the other hand, the oracles against the earlier nations are more
    specific, in that the specific rebellious acts are given. This shift
    emphasizes Judah's responsibility to God by contrasting their action with
    the high standards that God revealed to them. The other nations were judged
    on the basis of common sense and a popular sense of morality, while Judah
    had the statutes and laws which guided them.
    \autocite[115]{smith:2017}
\end{quote}

As Amos moves from the surrounding nations to Judah, the charges become more
specific. Judah has \enquote{rejected the law of the \textsc{Lord}} they have
\enquote{not kept his decrees.} Finally, we are told they \enquote{have been led
astray by false gods, / the gods their ancestors followed.}

\paragraph{Law} In our contemporary context, the word \enquote{law} carries with
it \enquote{related ideas of legalism, enforcement, reward, and punishment}
\autocite[54]{motyer:2011}, which is not quite how the original readers would
likely have understood it. For them, \enquote{law} would have meant
\enquote{\enquote{instruction}, with the related idea of a personal contact
between teacher and pupil.} It was a relational term, the type of instruction
that is given personally by a teachers to their students, not the type of
instruction that is handed down impersonally from a government to its citizens.
It is fitting, then, that Amos reminds the Israelites of their time in the
wilderness after the law had been given when God was with them and led them.
This law is not the law of a distant God who delivers a set of edicts one must
obey but the law of a loving father who instructs his children in the right way
to live for their own benefit.

\begin{quote}
    \enquote{Law} signifies not \enquote{legislation} (with its related ideas of
    legalism, enforcement, reward and punishment) but \enquote{instruction},
    with the related idea of a personal contact between teacher and pupil. The
    Lord's law implies His personal drawing near as teacher and the
    establishment of a personal relationship between Him and the one to whom He
    purposes to impart his truth.
    \autocite[54]{motyer:2011}
\end{quote}

\paragraph{Liars} The NIV translators made an interpretive decision when they
translated verse 4c as \enquote{they have been led astray by false gods, / the
gods their ancestors followed.} The original text simply says \enquote{lies,}
not \enquote{false gods.} Smith translates the phrase \enquote{but their liars
led them astray, / the ones after which their fathers went}
\autocite[106]{smith:2017}, noting that \enquote{the Old Greek translated
\enquote{their liars} with \enquote{their follies, idols} and adds
\enquote{which they have made} which is not in the Hebrew.}

\begin{quote}
    The final lines of 2:4 are sometimes interpreted to refer to Judah's worship
    of \enquote{false gods, or idols}. On the basis of this interpretation, a
    specific sin against Judah is identified, but it is a crime against Yahweh
    rather than a social sin as in the earlier oracles. But the root meaning of
    \textit{kzb} means \enquote{to lie}, and there is no other case in
    Deuteronomy (or the rest of Hebrew scripture) where this root clearly means
    \enquote{false gods, idols}. The prophets do describe the people as being
    \enquote{led astray} by the \enquote{lies} (\textit{kzb}) of the false
    prophets and leaders of the nation (Isa. 3:12; 9:15--16; 28:15; Mic. 3:5;
    Jer. 23:14, 32). These leaders deceived the people of Judah and caused
    perversion in the moral, civil, ceremonial, and personal behavior or
    attitudes of the people.
    \autocite[115-116]{smith:2017}
\end{quote}

There is clearly some debate about this and I have no skills, training or
experience in reading or understanding Hebrew, but it seems to me that
\enquote{liars,} referring to leaders who have led Israel astray is the
appropriate reading of these verses. The ESV, HCSB, KJV, NLT, and NRSV all
translate this part of verse 4 as \enquote{liars} with the NIV being the only
translation I looked at rendering it as \enquote{false gods.} I think what we
see here is a nation that has gradually abandoned God's plan through following
poor leaders.

\subsection{The Oracle Against Israel}

The oracle against Israel in verses 6--16 can be divided into three parts:

\begin{enumerate}
    \item The standard formula -- \enquote{For three sins of Israel, even for
        four, I will not relent} followed by a list of the sins that led to this
        judgement
    \item A reminder of what God has done -- they are God's special people, they
        have received many blessings and benefited from divine favour
    \item A description of the coming judgement -- it is inescapable
\end{enumerate}

The first six oracles against the nations surrounding Judah and Israel consisted
of accusations of immorality against other people. As I discussed in the post on
Amos 1, God is holding them to a general standard common to all people. In Amos
2, when the prophet's attention is turned to Judah and then to Israel, we see
that they are being held to account not only for sins against this general
standard everyone is expected to know and respect, but also for sins against the
more specific standard, the standard God revealed to them and for sins against
God. Judah, it says, \enquote{rejected the law of the \textsc{Lord}} (v. 4) and
Israel profaned God's holy name (v. 7).\footnote{%
    \enquote{By agreeing with Amos' accusations against Judah, the Israelite
    audience is admitting the validity of God's law as a standard for the
    judgment of his people. Judah knew God's will (and so did Israel), but she
    failed to follow him. Although failure to follow Yahweh's law is never
    specifically mentioned in the following oracle against Israel, most agree
    that the Israelite legal code is the backdrop for the accusation against
    Israel. Therefore, the establishment of its authority immediately before the
    Israelite oracle is significant. The specific sins of Israel (2:6--8) are
    contrary to the traditions which Judah and Israel received and the
    proclamation of the prophets (2:11). By getting Israel to support the
    judgment against Judah on the basis of the law of Yahweh, Amos has prepared
    the nation for God's word against the northern tribes.}
    \autocite[116]{smith:2017}
}

In verses 9--12, Amos retells some of the things God has done for his people:

\begin{itemize}
    \item He destroyed the Amorites who occupied the land before them (v. 9)
    \item He brought them out of Egypt (v. 10)
    \item He was with them, cared for them, led them through the wilderness (v.
        10)
    \item He gave them the land (v. 10)
    \item He never abandoned them, left them to their own devices, he raised up
        prophets and Nazirites (holy people) from among them (v. 11)
\end{itemize}

One of the key topics in Amos is, I think, the idea of privilege. God's people
had a unique privilege -- they had the law, they had instructions from God, they
knew what was right and what was wrong through more than just their conscience.
The list of what God has done for the Israelites was not simply a reminder that
he was good and worthy of their devotion and obedience. No, as we're told in
verse 6, God \enquote{will not relent.} This wasn't a call to repentance, this
was a notice of the coming judgment.

When Amos reminded the Israelites that the God whose name they had profaned was
the same God who destroyed the Amorites completely (from fruit to root), he was
reminding them that God is a God of justice, a powerful God who had destroyed a
powerful people before because of the very things the Israelites were now
themselves doing. In Amos 3:2 it says, \enquote{You only have I chosen / of all
the families of the earth; / therefore I will punish you for all your sins.} The
fact that they were God's chosen people didn't let them off the hook, God
wouldn't treat them more leniently than the Amorites. On the contrary, because
they had the privileges of being led by him through the wilderness, of having
prophets in their midst, of seeing holiness lived out in the lives of the
Nazirites, of possessing the law, they would be judged more harshly. They
wouldn't be let off the hook because God had been so good to them, they were all
the more aware of the seriousness of and responsible for their actions because
God had given them every assistance to live rightly and despite all this they
chose to live wrongly.

The list of things God had done for his people in verses 9--12 showed just how
completely they had rejected God's goodness and provision for them.
\enquote{Whatever their sins against the voice of conscience, this is no longer
in the forefront. They have passed the point of no return by disregarding and
silencing the voice of revelation.}
\autocite[109]{smith:2017}

\subsection{Israeli and Judah Together Under God's Judgement}

In 930 BC, Rehoboam began to reign and, as we learn in 1 Kings 12, due to
\enquote{the repressive measures his father, Solomon, had initiated to fund his
building projects}\autocite[22]{rusten:2005}

\begin{quote}
    Although Judah and Israel were once one nation, rivalry between the tribes
    existed before the reign of David (2 Sam. 2--3), and eventually the nation
    split after the death of Solomon (1 Kings 12). Bitterness and mistrust
    resulted in open war between the two brother nations (1 Kings 12:18--24;
    14:30; 15:7, 16--21, 32; 2 Kings 14:8--14). The introduction of the calves
    at Bethel and Dan, the acceptance of Baalism and other cultic innovations (1
    Kings 12:28--33) widened the gap between Judah and Israel. In many ways the
    Israelites in Samaria considered Judah a foreign nation, and the audience
    probably greeted Amos' condemnation of Judah with a good deal of pleasure
    and smug self-righteousness.
    \autocite[114]{smith:2017}
\end{quote}

The ten northern tribes that comprised the nation of Israel and the two southern
tribes that comprised the nation of Judah would be punished for the same reason:
because they were God's chosen people (with all the blessings and privileges
that came with it) and yet had still sinned, they would be punished.

Though Judah and Israel considered themselves separate nations, the list of
things God had done for the northern nation of Israel applies equally to the
southern tribe of Judah. They had all received the law, they had all seen God
destroy the Amorites, they had all been rescued by God from Egypt, they had all
been with God in the wilderness. As it says in Amos 3:1: \enquote{Hear this
word, people of Israel, the word the \textsc{Lord} has spoken against
you---against the whole family I brought up out of Egypt}.

\section{Theological Reflections}

\subsection{God Doesn't Play Favourites}

There is no distinction between the other nations and Israel and Judah. Amos
uses the same phrase to describe the repeated transgressions of Israel as he
used in the earlier oracles: \enquote{For three sins of Israel, / even for four,
I will not relent.} (2:4, 6)

Writing in his commentary on Amos, J. A. Motyer says this:

\begin{quote}
    In the same way, the appeal to conscience, to common humanity, underlying
    his review of the world is another move depriving Israel of any special
    ground or plea. Whatever makes Israel distinct among the nations, there is
    no distinction at this point, that the same moral rules operate inside as
    outside. Thus the noose tightens until, as we shall now see, the unique
    position granted by grace to the church of God, far from excusing or even
    ameliorating the offence, aggravates the situation so that Israel's fourth
    transgression is even less understandable or forgivable than that of the
    heathen who knew not God. The uniqueness of the church includes its unique
    peril.
    \autocite[50-51]{motyer:2011}
\end{quote}

This truth is directly applicable to the church today. We have been redeemed
through the blood of Christ. This is our rescue-from-Egypt story. But the idea
that we can sit back and live however we like is totally incompatible with
everything the Bible teaches through both the Old and New Testaments. It is an
age-old tension -- we are the rescued and redeemed people of God but that
doesn't lessen our responsibility, it increases it.

\begin{quote}
    \enquote{I will not vacillate about it} demonstrates Yahweh's determination
    to hold Judah responsible for her guilt. His wrath will not cease; the
    decree will be carried out; Judah will be punished for her many rebellious
    acts. Any privilege which Judah enjoyed through her special relationship to
    God is ignored. The possibility of excuse is not even mentioned; instead the
    advantage of Judah increases the weight of her guilt.
    \autocite[114-115]{smith:2017}
\end{quote}

In Luke 12:42ff, Jesus contrasts a \enquote{faithful and wise manager} with a
\enquote{servant who knows the master's will and \dots does not do what the
master wants.} Jesus said \enquote{it will be good for [the faithful and wise]
servant \dots when [the master] returns} but the master \enquote{will cut [the
second sedrvant] to pieces and assign him a place with unbelievers.}

\begin{quote}
    \enquote{The servant who knows the master's will and does not get ready or
    does not do what the master wants will be beaten with many blows. But the
    one who does not know and does things deserving punishment will be beaten
    with few blows. From everyone who has been given much, much will be
    demanded; and from the one who has been entrusted with much, much more will
    be asked.} (Luke 12:47--48)
\end{quote}

Like Amos, Jesus tells us that knowing the truth isn't enough. On the contrary,
those who know what the master wants and don't do it will receive a harsher
punishment. There is no favouritism with God. But what does this mean? Perhaps
there is more to it than this, but I think at least this means that growing up
in a Christian family or learning about the Christian faith without actually
submitting to Jesus as your Lord and God as your master isn't going to help one
bit. There are plenty of people who call themselves Christians for whom these
must be terrifying verses indeed. For me, they give me pause for thought and I'd
like to continue to try to understand exactly what they mean.

\subsection{God's People Are Vulnerable to Lies}

In Amos, we see God's chosen people falling into the sins of the Amorites. We
see a gradual corruption through tradition: \enquote{because they have been led
astray by false gods, the gods their ancestors followed} (v. 4). It happened to
God's people in the Old Testament and we would, I think, be fools to think that
it cannot happen to the church today.

Motyer made two relevant points on this topic:

\begin{enumerate}
    \item God's truth (as revealed to us in the Bible) is \enquote{the safeguard
        of the church} just as much as \enquote{the church is called by God to
        safeguard} it. We are to look to the Bible, not to man (no matter how
        good their theological pedigree) for the truth about how we should live.
        It's important that we don't let the idea that we are the guardians of
        the truth drown out the Bible's ability to speak to us and correct our
        understanding no matter how much tradition is on the line. The Bible
        keeps the Christian safe from bad leadership in the church just as much
        as the church proclaims the truths of the Bible.

        \begin{quote}
            Here is a very largely forgotten and yet most vital principle. It is
            certainly the case that the church is called by God to safeguard,
            publicize and transmit His truth (\textit{e.g.} 2 Tim. 1:13, 14;
            2:1, 2); but it is equally the case that the truth is the safeguard
            of the church, both in the corporate sense of preserving the whole
            body and in the individual sense of guarding, defending and keeping
            each member. The life which walks int the truth is impregnable
            (\textit{cf.} Jn. 8:31, 32, 34--36).
            \autocite[54]{motyer:2011}
        \end{quote}

    \item The Bible is not the only source of truth. In Proverbs 6:6, we're
        told, \enquote{Go to the ant, you sluggard; / consider its ways and be
        wise!} We can look at the world God has created and we can listen to the
        people God has created, the people who bear his image, to learn truth.
        But the instruction God has revealed to us in the Bible must be our
        supreme source of truth. That's the thing we are to base our lives on.

        \begin{quote}
            Now Amos is not saying that man is incapable of ever unearthing
            truth. He is not even saying that unless a thing can be verified by
            Scripture it must be false. He is saying this: when anything other
            than the Word of God is given the supreme place, so that we base our
            lives upon it and guide our lives by it, then it becomes a lie and a
            source of lies. To put it at its best---\textit{i.e.} when
            \enquote{tradition} expresses a truth and not an error---when a part
            of the truth is taken to be the whole of the truth it becomes an
            untruth, simply because it is stretched beyond its limit, and
            becomes warped and altered. And when such a misshapen
            \enquote{truth} is taken as the guide for life, what can it do but
            mislead?
            \autocite[54]{motyer:2011}
        \end{quote}
\end{enumerate}

I think the words of this oracle, that \enquote{They have rejected the
instruction of the \textsc{Lord}, / refusing to obey his decrees. / They have
been led astray by the same lies / that deceived their ancestors} (Amos 2:4b--c,
NLT) speak directly into our context today. We need to make sure we never forget
God's instruction (and the type of instruction it is, i.e., \textit{relational}
instruction) and fall victim to the same lies humans have been deceived by for
as long as there have been humans.

\section{We Cannot Save Ourselves}

Verses 13--15 expand on the coming judgement. Amos corrects any mistaken idea in
the people that there would be hope for escape. God will crush them and nobody
will escape, not the swift, the strong, the warrior, the teacher, the archer,
the fleet-footed soldier, the horseman. The bravest warriors will be humiliated
-- they will flee naked on the coming day of judgement.

Niehaus sees hope in verses 13ff and certainly this is in line with God's
character, but I'm not sure it can be tied back to Amos' message here.

\begin{quote}
    In the midst of all this, however, there is ultimate hope. For it is the
    Lord who chastises them. And the Lord is, after all, \enquote{God the Hero.}
    \enquote{The Lord disciplines those whom he loves, / and chastises every
    child whom he accepts} (Heb. 12:6, quoting Prov. 3:12). And he must win the
    battle.
    \autocite[372]{mccomiskey:2009}
\end{quote}

Those we rely on for help, those to whom we look to save us will not be able to
do so. Against the punishment of God, all human effort fails.

\begin{quote}
    It illustrates, along with verses 14--15, the principle there stated, that
    none of the able, skilled, or mighty will be able to do what is expected of
    them when the judgement of Yahweh comes. Indeed, they will not even be able
    to deliver themselves, not to mention making an effective stand against the
    foe.
    \autocite[373]{mccomiskey:2009}
\end{quote}

\section*{Postscript}

One of the things I've noticed as I've been studying Amos is how much I don't
yet understand about God's story. The Old and New Testaments must be understood
as a whole, but I find so many unanswered questions about passages in the New
Testament getting in the way of me understanding the Old. For example, what does
it mean that there is no condemnation for believers (cf. Romans 8:1)? Does it
make Amos' warnings irrelevant to the Christian? Are they purely of historical
value? Do they reveal the will and heart of God? I haven't given a lot of
thought to Romans 8:1 but it does seem to be related to Amos' words that God's
anger would rage against the Israelites and nothing could turn it away. I'm sure
I'll understand more things as I continue to engage with God's word; I'm also
sure I'll never get to the end, I'll never know everything. Perhaps when I study
Romans, I'll find that my ideas about how to interpret Amos were wrong.

Another example of this is the New Testament teaching about obeying secular
authorities in Romans 13 (\enquote{Let everyone be subject to the governing
authorities, for there is no authority except that which God has established.}
[Romans 13:1]) and 1 Peter 2 (\enquote{Submit yourselves for the Lord's sake to
every human authority: whether to the emperor, as the supreme authority, or to
governors, who are sent by him to punish those who do wrong and to commend those
who do right.} [1 Peter 2:13]). In Amos 2:4, the Israelites are accused of
following poor leaders. Perhaps there's some extra context since the nation of
Israel and its leaders were chosen by God in a way that our secular nations are
not? I think it quite likely that it's a warning against following ungodly
leaders in our churches.

\newpage

\thispagestyle{first}

\mbox{}\vspace{1in} % Extra margin

\printbibliography[title={BIBLIOGRAPHY}]

\end{document}

